\documentclass{article}
\usepackage{amsmath}
\usepackage{amssymb}
\usepackage{siunitx}
\usepackage{chemformula}

\begin{document}

\section{原子结构}
\subsection{原子光谱与玻尔理论}

\subsubsection{原子结构理论发展历程}
\begin{itemize}
    \item \textbf{1803年 Dalton 原子论}：提出原子不可再分概念
    \item \textbf{1897年 Thomson 模型}：阴极射线实验测得电子荷质比 $\frac{e}{m_e} = \SI{1.758e11}{\coulomb\per\kilo\gram}$
    \item \textbf{1909年 Millikan 油滴实验}：测定电子电量 $e = \SI{1.592e-19}{\coulomb}$
    \item \textbf{1911年 Rutherford 模型}：$\alpha$ 粒子散射实验提出核式结构
\end{itemize}

\subsubsection{氢原子光谱特征}
\begin{itemize}
    \item \textbf{Rydberg 公式}：
        $\tilde{\nu} = R\left(\frac{1}{n_1^2} - \frac{1}{n_2^2}\right)$
    \item \textbf{里德堡常数}：$R = \SI{1.097e7}{\per\meter}$
    \item \textbf{光谱类型}：线状光谱，不连续特征
\end{itemize}

\subsubsection{玻尔理论核心内容}
\begin{itemize}
    \item \textbf{普朗克量子假说}：$E = h\nu$，$h = \SI{6.626e-34}{\joule\second}$
    \item \textbf{轨道量子化条件}：
        \begin{align*}
        r &= \SI{52.9}{\pico\meter} \times \frac{n^2}{Z} \\
        E &= -\SI{21.8e-19}{\joule} \times \frac{Z^2}{n^2}
        \end{align*}
    \item \textbf{能级跃迁公式}：$\Delta E = E_2 - E_1 = h\nu$
\end{itemize}

\subsubsection{微观粒子波粒二象性}
\begin{itemize}
    \item \textbf{德布罗意关系}：$\lambda = \frac{h}{p} = \frac{h}{mv}$
    \item \textbf{海森堡不确定性原理}：$\Delta x \cdot \Delta p \geq \frac{h}{4\pi}$
    \item \textbf{实验验证}：戴维森-革末电子衍射实验
\end{itemize}

\subsection{电子在核外运动状态的描述}

\subsubsection{薛定谔方程基本形式}
\[
\frac{\partial^2 \psi}{\partial x^2} + \frac{\partial^2 \psi}{\partial y^2} + \frac{\partial^2 \psi}{\partial z^2} + \frac{8\pi^2 m}{h^2}(E - V)\psi = 0
\]

\subsubsection{四个量子数及其物理意义}
\begin{itemize}
    \item \textbf{主量子数 $n$}：决定电子层和主要能量
    \item \textbf{角量子数 $l$}：决定轨道形状和角动量
        $|M| = \frac{h}{2\pi}\sqrt{l(l+1)}$
    \item \textbf{磁量子数 $m$}：决定空间取向，$m = 0,\pm1,\dots,\pm l$
    \item \textbf{自旋量子数 $m_s$}：$m_s = \pm\frac{1}{2}$
\end{itemize}

\subsubsection{原子轨道与电子云}
\begin{itemize}
    \item \textbf{概率密度}：$|\psi|^2$ 表示电子出现概率密度
    \item \textbf{径向分布函数}：$D(r) = 4\pi r^2 |\psi|^2$
    \item \textbf{玻尔半径}：$a_0 = \SI{52.9}{\pico\meter}$
\end{itemize}

\subsubsection{电子层结构规律}
\begin{center}
\begin{tabular}{ccccc}
\toprule
$n$ & 亚层 & 轨道数 & 电子层容量 & 最大电子数 \\
\midrule
1 & $1s$ & 1 & $n^2$ & $2n^2$ \\
2 & $2s,2p$ & 4 & 4 & 8 \\
3 & $3s,3p,3d$ & 9 & 9 & 18 \\
4 & $4s,4p,4d,4f$ & 16 & 16 & 32 \\
\bottomrule
\end{tabular}
\end{center}

\subsection{多电子原子结构}

\subsubsection{核外电子排布三原则}
\begin{itemize}
    \item \textbf{泡利不相容原理}：同一轨道最多两个电子，自旋相反
    \item \textbf{能量最低原理}：电子优先占据能量最低轨道
    \item \textbf{洪特规则}：简并轨道电子分占且自旋平行
\end{itemize}

\subsubsection{能级组与周期关系}
\begin{center}
\begin{tabular}{cccc}
\toprule
能级组 & 包含轨道 & 电子容量 & 对应周期 \\
\midrule
1 & $1s$ & 2 & 第一周期 \\
2 & $2s,2p$ & 8 & 第二周期 \\
3 & $3s,3p$ & 8 & 第三周期 \\
4 & $4s,3d,4p$ & 18 & 第四周期 \\
\bottomrule
\end{tabular}
\end{center}

\subsubsection{元素周期表分区}
\begin{center}
\begin{tabular}{ccc}
\toprule
分区 & 族 & 价电子构型 \\
\midrule
$s$区 & \text{I A}、\text{II A} & $ns^{1-2}$ \\
$p$区 & \text{III A}$\sim$\text{VII A}、0族 & $ns^2 np^{1-6}$ \\
$d$区 & \text{III B}$\sim$\text{VII B}、\text{VIII}族 & $(n-1)d^{1-8} ns^{1-2}$ \\
$ds$区 & \text{I B}、\text{II B} & $(n-1)d^{10} ns^{1-2}$ \\
$f$区 & 镧系、锕系 & $(n-2)f^{0-14} (n-1)d^{0-1} ns^2$ \\
\bottomrule
\end{tabular}
\end{center}

\subsection{原子的基本性质}

\subsubsection{有效核电荷计算}
\begin{itemize}
    \item \textbf{屏蔽效应公式}：$Z' = Z - \sum\sigma$
    \item \textbf{计算示例}：
        \begin{itemize}
            \item \ch{K}：$Z' = 19 - 16.8 = 2.2$
            \item \ch{Ca}：$Z' = 20 - 17.15 = 2.85$
            \item \ch{Sc}：$Z' = 21 - 17.9 = 3.00$
        \end{itemize}
\end{itemize}

\subsubsection{原子半径变化规律}
\begin{itemize}
    \item \textbf{定义方式}：共价半径、金属半径、范德华半径
    \item \textbf{周期变化}：同一周期从左到右半径减小
    \item \textbf{族变化}：同一主族从上到下半径增大
    \item \textbf{镧系收缩}：$5\rightarrow6$ 周期半径相近
\end{itemize}

\subsubsection{电离能与电负性}
\begin{itemize}
    \item \textbf{第一电离能}：$\ch{A(g) -> A^+(g) + e^-}$
    \item \textbf{电负性标度}：Pauling标度，\ch{F}最大($3.98$)，\ch{Cs}最小($0.79$)
    \item \textbf{金属性判据}：电负性通常小于$2.0$为金属
\end{itemize}

\end{document}